\documentclass[a4paper]{article}

% 文章排版
\usepackage{indentfirst}
\setlength{\parindent}{0em}
\usepackage{autobreak}
\allowdisplaybreaks
\usepackage{parskip}
\usepackage{geometry}
\geometry{top=3cm,bottom=3cm,left=2cm,right=2cm}
\usepackage{anyfontsize}

% 数学物理宏包
\usepackage{amsmath,amssymb}
\usepackage{tabularray}
\UseTblrLibrary{amsmath}
\usepackage{physics}
\renewcommand\div\divisionsymbol
\DeclareMathOperator{\diag}{diag}
\DeclareMathOperator{\artanh}{artanh}
\DeclareMathOperator{\arsech}{arsech}
\newcommand{\trans}{\text{\textsf{T}}}

\usepackage{fontspec}
\setmainfont{Lucida Bright OT}
\usepackage{unicode-math}
\unimathsetup{mathrm=sym,mathbf=sym,mathsf=sym,math-style=ISO,bold-style=ISO}
\setmathfont[StylisticSet=4,sans-style=italic]{Lucida Bright Math OT}
\setmathfont[range={\mathscr,\mathbfscr}]{Lucida Bright Math OT}

\usepackage[fontset=none]{ctex}
\setCJKmainfont{FZShuSong-Z01}[BoldFont=Source Han Sans CN,ItalicFont=FZKai-Z03]
\setCJKsansfont{Source Han Sans CN}
\usepackage{tikz}
\usetikzlibrary{arrows.meta}
\usetikzlibrary{decorations.pathmorphing}
\usetikzlibrary{decorations.pathreplacing}

% 符号重定义
\AtBeginDocument{
    \let\uglyepsilon\epsilon
    \renewcommand\epsilon\varepsilon
    \let\uglyleq\leq
    \renewcommand\leq\leqslant
    \let\uglygeq\geq
    \renewcommand\geq\geqslant
}


% 题目
\title{\textbf{广义相对论\ \ \ \ 第四次作业}}
\author{张植竣 1900011604}
\date{2023年6月18日}

\begin{document}

\maketitle

\section*{第十章}

\begin{enumerate}
    \item[9.]
    在Boyer-Lindquist坐标下，Kerr时空几何的线元可写作：（$\mu = GM$，下同）
    \[
        \dd{s}^2 = -\dd{t}^2 + \frac{\rho^2}{\Delta} \dd{r}^2 + \rho^2 \dd{\theta}^2 + (r^2 + a^2)\sin^2{\theta} \dd{\phi}^2 + \frac{2\mu r}{\rho^2} \qty(a\sin^2{\theta}\dd{\phi} - \dd{t})^2
    \]
    或者写作：
    \[
        \dd{s}^2 = -\frac{\rho^2 \Delta}{\Sigma^2} \dd{t}^2 + \frac{\rho^2}{\Delta} \dd{r}^2 + \rho^2 \dd{\theta}^2 + \frac{\Sigma^2 \sin^2{\theta}}{\rho^2} \qty(\dd{\phi} - \omega \dd{t})^2
    \]
    其中：
    \[
        \rho^2 = r^2 + a^2 \cos\theta,\quad \Delta = r^2 - 2\mu r + a^2,\quad \omega = \frac{2\mu ra}{\Sigma^2}
    \]
    以及：
    \[
        \Sigma^2 = (r^2 + a^2)^2 - a^2 \Delta \sin^2{\theta}
    \]
    可以证明，它同时也等于：
    \[
        \Sigma^2 = (r^2 + a^2) \rho^2 + 2\mu r a^2 \sin^2{\theta}
    \]
    这个关系在两种写法的转换中非常有用。
    \begin{itemize}
        \item 对于由$t = \mathrm{const.}$和事件视界$r = r_\pm$定义的曲面：
        \[
            r_\pm = \mu \pm \sqrt{\mu^2 - a^2},\quad \Delta_\pm = 0,\quad \Sigma_\pm = (r_\pm^2 + a^2)^2 = (2\mu r_\pm)^2
        \]
        度规线元为：
        \[
            \dd{s}^2 = \rho_\pm^2 \dd{\theta}^2 + \frac{\Sigma_\pm^2}{\rho_\pm^2} \sin^2{\theta} \dd{\phi}^2 = \rho_\pm^2 \dd{\theta}^2 + \qty(\frac{2\mu r_\pm}{\rho_\pm})^2 \sin^2{\theta} \dd{\phi}^2
        \]
        其中：
        \[
            \rho_\pm^2 = r_\pm^2 + a^2\cos^2{\theta}
        \]
        绕两极的大圆$\dd{\phi} = 0$，则周长为：
        \[
            L_1 = 2\int_0^\pi \sqrt{g_{\theta\theta}(r_\pm)} \dd{\theta} = 2\int_0^\pi \rho_\pm \dd{\theta} = 2\int_0^\pi \sqrt{r_\pm^2 + a^2\cos^2{\theta}} \dd{\theta}
        \]
        绕赤道的大圆$\theta = \dfrac{\pi}{2}$，$\rho_\pm^2 = r_\pm^2 + a^2\cos^2{\theta} = r_\pm^2$，即$\rho_\pm = r_\pm$，则周长为：
        \[
            L_2 = \int_0^{2\pi} \sqrt{g_{\phi\phi}\qty(r_\pm, \theta=\frac{\pi}{2})} \dd{\phi} = \int_0^{2\pi} \qty(\frac{2\mu r_\pm}{\rho_\pm}) \sin{\theta} \dd{\phi} = \int_0^{2\pi} \qty(\frac{2\mu r_\pm}{r_\pm}) \dd{\phi} = 4\pi\mu
        \]
        因此我们需要证明：
        \[
            2\int_0^\pi \sqrt{r_\pm^2 + a^2\cos^2{\theta}} \dd{\theta} < 4\pi\mu
        \]
        即：
        \[
            \int_0^\pi \sqrt{r_\pm^2 + a^2\cos^2{\theta}} \dd{\theta} < 2\pi\mu
        \]
        由于$0 \leqslant \cos^2{\theta} \leqslant 1$，则：（注意这个不等式是取不到等号的）
        \[
            \int_0^\pi \sqrt{r_\pm^2 + a^2\cos^2{\theta}} \dd{\theta} < \int_0^\pi \sqrt{r_\pm^2 + a^2} \dd{\theta}
        \]
        而：
        \[
            \int_0^\pi \sqrt{r_\pm^2 + a^2} \dd{\theta} \leqslant \int_0^\pi \sqrt{r_+ + a^2} \dd{\theta} = 2\pi \sqrt{r_+ + a^2}
        \]
        又因为：
        \[
            r_+ = \mu + \sqrt{\mu^2 - a^2} \leqslant 2\mu
        \]
        且上式只有在$a = J/M = 0$时取等号，此时退化为Schwarschild黑洞。因此对于Kerr黑洞，一定有：
        \[
            \int_0^\pi \sqrt{r_\pm^2 + a^2\cos^2{\theta}} \dd{\theta} < 2\pi\mu
        \]
        即对于由$t = \mathrm{const.}$和$r = r_\pm$定义的曲面，绕两极的圆周长小于沿赤道的圆周长。

        \item 对于由$t = \mathrm{const.}$和无穷大红移面$r = r_{S\pm}$定义的曲面：
        \[
            g_{tt} = -1 + \frac{2\mu r}{\rho^2} = 0 \implies r_{S\pm} = \mu \pm \sqrt{\mu^2 - a^2\cos^2\theta}
        \]
        则：
        \[
            \rho_{S\pm}^2 = 2\mu r_{S\pm},\quad \Delta_{S\pm} = a^2 \sin^2{\theta}
        \]
        注意，此时$r$是$\theta$的函数，$\dd{r} \neq 0$，度规线元为：
        \begin{align*}
            \dd{s}^2
            &= \frac{\rho_{S\pm}^2}{\Delta_{S\pm}} \dd{r}^2 + \rho_{S\pm}^2 \dd{\theta}^2 + \frac{\Sigma_{S\pm}^2}{\rho_{S\pm}^2} \sin^2{\theta} \dd{\phi}^2 \\
            &= \frac{2\mu r_{S\pm}}{a^2 \sin^2{\theta}} \dd{r}^2 + 2\mu r_{S\pm} \dd{\theta}^2 + \frac{(r_{S\pm}^2 + a^2) \rho_{S\pm}^2 + 2\mu r_{S\pm} a^2 \sin^2{\theta}}{\rho_{S\pm}^2} \sin^2{\theta} \dd{\phi}^2 \\
            &= \frac{2\mu r_{S\pm}}{a^2 \sin^2{\theta}} \dd{r}^2 + 2\mu r_{S\pm} \dd{\theta}^2 + \frac{(r_{S\pm}^2 + a^2) \rho_{S\pm}^2 + \rho_{S\pm}^2 a^2 \sin^2{\theta}}{\rho_{S\pm}^2} \sin^2{\theta} \dd{\phi}^2 \\
            &= \frac{2\mu r_{S\pm}}{a^2 \sin^2{\theta}} \dd{r}^2 + 2\mu r_{S\pm} \dd{\theta}^2 + \qty(r_{S\pm}^2 + a^2 + a^2 \sin^2{\theta}) \sin^2{\theta} \dd{\phi}^2 \\
            &= \frac{2\mu r_{S\pm}}{a^2 \sin^2{\theta}} \dd{r}^2 + 2\mu r_{S\pm} \dd{\theta}^2 + \qty(\rho_{S\pm}^2 + 2a^2 \sin^2{\theta}) \sin^2{\theta} \dd{\phi}^2 \\
            &= \frac{2\mu r_{S\pm}}{a^2 \sin^2{\theta}} \dd{r}^2 + 2\mu r_{S\pm} \dd{\theta}^2 + 2\qty(\mu r_{S\pm} + a^2 \sin^2{\theta}) \sin^2{\theta} \dd{\phi}^2
        \end{align*}
        绕两极的大圆$\dd{\phi} = 0$，则周长为：
        \[
            L_1 = \int_L \dd{s} = \int_L \sqrt{\frac{2\mu r_{S\pm}}{a^2 \sin^2{\theta}} \dd{r}^2 + 2\mu r_{S\pm} \dd{\theta}^2}
        \]
        但我们只能选择一个积分参数，选择$\theta$作为积分参数，则有：
        \[
            L_1 = 2\int_0^\pi \sqrt{\frac{2\mu r_{S\pm}}{a^2 \sin^2{\theta}} \qty(\dv{r_{S\pm}}{\theta})^2 + 2\mu r_{S\pm}} \dd{\theta}
        \]
        其中：
        \[
            \dv{r_{S\pm}}{\theta} = \dv{\theta} \qty(\mu \pm \sqrt{\mu^2 - a^2\cos^2\theta}) = \pm \frac{a^2 \sin{\theta} \cos{\theta}}{\sqrt{\mu^2 - a^2 \cos^2{\theta}}} = \frac{a^2 \sin{\theta} \cos{\theta}}{r_{S\pm} - \mu}
        \]
        这一步也可以通过对$2\mu r_{S\pm} = \rho_{S\pm}^2 = r_{S\pm}^2 + a^2\cos^2{\theta}$两边求导得到。

        于是有：
        \[
            \frac{2\mu r_{S\pm}}{a^2 \sin^2{\theta}} \qty(\dv{r_{S\pm}}{\theta})^2
            = \frac{2\mu r_{S\pm}}{a^2 \sin^2{\theta}} \cdot \frac{a^4 \sin^2{\theta} \cos^2{\theta}}{(r_{S\pm} - \mu)^2}
            = \frac{2\mu r_{S\pm}}{(r_{S\pm} - \mu)^2} a^2 \cos^2{\theta}
        \]
        注意$r_{S\pm} = \mu \pm \sqrt{\mu^2 - a^2\cos^2\theta}$，则$(r_{S\pm} - \mu)^2 = \mu^2 - a^2\cos^2\theta$，则：
        \begin{align*}
            &\quad \frac{2\mu r_{S\pm}}{a^2 \sin^2{\theta}} \qty(\dv{r_{S\pm}}{\theta})^2 + 2\mu r_{S\pm} \\
            &= \frac{2\mu r_{S\pm}}{(r_{S\pm} - \mu)^2} a^2 \cos^2{\theta} + 2\mu r_{S\pm} \\
            &= \frac{2\mu r_{S\pm}}{(r_{S\pm} - \mu)^2} \qty[a^2 \cos^2{\theta} + (r_{S\pm} - \mu)^2] \\
            &= \frac{2\mu r_{S\pm}}{(r_{S\pm} - \mu)^2} \qty[a^2 \cos^2{\theta} + \mu^2 - a^2 \cos^2{\theta}] \\
            &= \frac{2\mu r_{S\pm}}{(r_{S\pm} - \mu)^2} \mu^2 \\
            &= \frac{2\mu^3 r_{S\pm}}{\mu^2 - a^2\cos^2\theta}
        \end{align*}
        因此：
        \[
            L_1 = 2\int_0^\pi \sqrt{\frac{2\mu^3 r_{S\pm}}{\mu^2 - a^2\cos^2\theta}} \dd{\theta} = 2\sqrt{2\mu^3} \int_0^\pi \sqrt{\frac{\mu \pm \sqrt{\mu^2 - a^2\cos^2\theta}}{\mu^2 - a^2\cos^2\theta}} \dd{\theta}
        \]
        绕赤道的大圆$\theta = \dfrac{\pi}{2}$，$r_{S\pm} = \mu \pm \mu$，则周长为：
        \[
            L_2 = \int_0^{2\pi} \sqrt{g_{\phi\phi}\qty(r_{S\pm}, \theta=\frac{\pi}{2})} \dd{\phi} = \int_0^{2\pi} \sqrt{2\qty(\mu r_{S\pm} + a^2)} \dd{\phi} = 2\pi \sqrt{2\qty(\mu r_{S\pm} + a^2)}
        \]
        代入$r_{S\pm} = \mu \pm \mu$，则有：
        \[
            L_{2+} = 2\pi \sqrt{2\qty(2\mu^2 + a^2)},\quad L_{2-} = 2\sqrt{2}\pi a
        \]
        \textcolor{red}{（我不会证明$L_1 < L_2$，考完试再说吧）}
    \end{itemize}

    \item[10.]
    对于外视界$r_+ = \mu + \sqrt{\mu^2 - a^2}$，$r$与$\theta$、$\phi$无关，则Kerr时空外视界面的度规分量为：
    \[
        \dd{s}^2 = \rho_\pm^2 \dd{\theta}^2 + \frac{\Sigma_\pm^2}{\rho_\pm^2} \sin^2{\theta} \dd{\phi}^2 = \rho_\pm^2 \dd{\theta}^2 + \qty(\frac{2\mu r_\pm}{\rho_\pm})^2 \sin^2{\theta} \dd{\phi}^2
    \]
    面元为：
    \[
        \dd{A} = \sqrt{|g_{\mu\nu}|} \dd{x^\mu} \dd{x^\nu} = 2\mu r_+ \sin{\theta} \dd{\theta} \dd{\phi}
    \]
    因此，外视界面积为：
    \[
        A = \int_0^{2\pi} \dd{\phi} \int_0^\pi 2\mu r_+ \sin{\theta} \dd{\theta} = 4\pi (2\mu r_+) = 8\pi \mu \qty(\mu + \sqrt{\mu^2 - a^2})
    \]
    对于任意两个分别具有质量$M_1$和$M_2$、角动量$J_1$和$J_2$的Kerr黑洞，视界面积分别为：
    \[
        A_1 = 8\pi \mu_1 \qty(\mu_1 + \sqrt{\mu_1^2 - a_1^2}),\quad \mu_1 = GM_1,\quad a_1 = \frac{J_1}{M_1}
    \]
    \[
        A_2 = 8\pi \mu_2 \qty(\mu_2 + \sqrt{\mu_2^2 - a_2^2}),\quad \mu_2 = GM_2,\quad a_2 = \frac{J_2}{M_2}
    \]
    而质量为$M_1 + M_2$、角动量为$J_1 + J_2$的Kerr黑洞的视界面积为：
    \[
        A_{12} = 8\pi \mu_{12} \qty(\mu_{12} + \sqrt{\mu_{12}^2 - a_{12}^2}),\quad \mu_{12} = G(M_1 + M_2),\quad a_{12} = \frac{J_1 + J_2}{M_1 + M_2} = \frac{\mu_1 a_1 + \mu_2 a_2}{\mu_1 + \mu_2}
    \]
    下面证明$A_1 + A_2 < A_{12}$，即：
    \[
        8\pi \mu_1 \qty(\mu_1 + \sqrt{\mu_1^2 - a_1^2}) + 8\pi \mu_2 \qty(\mu_2 + \sqrt{\mu_2^2 - a_2^2}) < 8\pi \mu_{12} \qty(\mu_{12} + \sqrt{\mu_{12}^2 - a_{12}^2})
    \]
    这等价于：
    \[
        \mu_1^2 + \mu_1 \sqrt{\mu_1^2 - a_1^2} + \mu_2^2 + \mu_2 \sqrt{\mu_2^2 - a_2^2} < (\mu_1 + \mu_2)^2 + (\mu_1 + \mu_2) \sqrt{(\mu_1 + \mu_2)^2 - \frac{(\mu_1 a_1 + \mu_2 a_2)^2}{(\mu_1 + \mu_2)^2}}
    \]
    化简得：
    \[
        \mu_1^2 + \mu_1 \sqrt{\mu_1^2 - a_1^2} + \mu_2^2 + \mu_2 \sqrt{\mu_2^2 - a_2^2} < (\mu_1 + \mu_2)^2 + \sqrt{(\mu_1 + \mu_2)^4 - (\mu_1 a_1 + \mu_2 a_2)^2}
    \]
    这个不等式可以拆分成两部分：
    \begin{equation}
        \mu_1^2 + \mu_2^2 < (\mu_1 + \mu_2)^2
    \end{equation}
    以及：
    \begin{equation}
        \mu_1 \sqrt{\mu_1^2 - a_1^2} + \mu_2 \sqrt{\mu_2^2 - a_2^2} < \sqrt{(\mu_1 + \mu_2)^4 - (\mu_1 a_1 + \mu_2 a_2)^2}
    \end{equation}
    显然，不等式(1)成立。下面证明不等式(2)也成立：

    (2)式可以写为：
    \[
        \sqrt{\mu_1^4 - \mu_1^2 a_1^2} + \sqrt{\mu_2^4 - \mu_2^2 a_2^2} < \sqrt{(\mu_1 + \mu_2)^4 - (\mu_1 a_1 + \mu_2 a_2)^2}
    \]
    不等式右边可以做如下放缩：
    \[
        (\mu_1 + \mu_2)^4 > (\mu_1^2 + \mu_2^2)^2
    \]
    该式成立的原因是其两边开平方后：
    \[
        (\mu_1 + \mu_2)^2 = \mu_1^2 + \mu_2^2 + 2\mu_1 \mu_2 > \mu_1^2 + \mu_2^2
    \]
    在$\mu > 0$时成立。

    那么只要证明放缩后的不等式(2)：
    \begin{equation}
        \sqrt{\mu_1^4 - \mu_1^2 a_1^2} + \sqrt{\mu_2^4 - \mu_2^2 a_2^2} \leqslant \sqrt{(\mu_1^2 + \mu_2^2)^2 - (\mu_1 a_1 + \mu_2 a_2)^2}
    \end{equation}
    成立，即可证明原不等式(2)成立，证明如下：

    记$x = \mu^2$，$y = \mu a$，则(3)式可写为：
    \[
        \sqrt{x_1^2 - y_1^2} + \sqrt{x_2^2 - y_2^2} \leqslant \sqrt{(x_1 + x_2)^2 - (y_1 + y_2)^2}
    \]
    两边平方得：
    \[
        \mathrm{LHS}^2 = (x_1^2 - y_1^2) + (x_2^2 - y_2^2) + 2\sqrt{(x_1^2 - y_1^2)(x_2^2 - y_2^2)}
    \]
    \[
        \mathrm{RHS}^2 = (x_1 + x_2)^2 - (y_1 + y_2)^2 = (x_1^2 - y_1^2) + (x_2^2 - y_2^2) + 2(x_1 x_2 - y_1 y_2)
    \]
    消去两边的公共项，得到：
    \[
        \sqrt{(x_1^2 - y_1^2)(x_2^2 - y_2^2)} \leqslant x_1 x_2 - y_1 y_2
    \]
    两边平方得：
    \[
        (x_1^2 - y_1^2)(x_2^2 - y_2^2) < (x_1 x_2 - y_1 y_2)^2
    \]
    展开得：
    \[
        \mathrm{LHS} = x_1^2 x_2^2 - x_1^2 y_2^2 - y_1^2 x_2^2 + y_1^2 y_2^2
    \]
    \[
        \mathrm{RHS} = x_1^2 x_2^2 - 2x_1 x_2 y_1 y_2 + y_1^2 y_2^2
    \]
    消去两边的公共项，得到：
    \[
        -x_1^2 y_2^2 - y_1^2 x_2^2 \leqslant -2x_1 x_2 y_1 y_2
    \]
    即：
    \[
        x_1^2 y_2^2 + y_1^2 x_2^2 \geqslant 2x_1 x_2 y_1 y_2
    \]
    这恰好可以转化为一个完全平方式：
    \[
        (x_1 y_2 - x_2 y_1)^2 \geqslant 0
    \]
    其中：
    \[
        x_1 y_2 - x_2 y_1 = \mu_1^2 \mu_2 a_2 - \mu_2^2 \mu_1 a_1 = \mu_1 \mu_2 (\mu_1 a_2 - \mu_2 a_1)
    \]
    只有在$\dfrac{a_1}{\mu_1} = \dfrac{a_2}{\mu_2}$时取等号。因此不等式(3)成立。进而不等式(2)成立，且取不到等号。

    综上所述，不等式(1)和(2)均成立，且均取不到等号，因此原不等式严格成立，即：
    \[
        A_1 + A_2 < A_{12}
    \]
\end{enumerate}

\section*{第十一章}

\begin{enumerate}
    \item[8.(a)]
    对于一颗能量密度为$\rho(r, t)$的球对称恒星，能动张量的00分量就是其能量密度：
    \[
        T^{00} = \rho(r, t)
    \]
    四极矩的定义为：
    \[
        q_{ij}(t) = \iiint_V r^i r^j T^{00}(t, \mathbf{r}) \dd^3{r}
    \]
    因为球对称性：
    \[
        q_{xx} = q_{yy} = q_{zz} = \iiint_V x^2 \rho(r, t) \dd^3{r}
    \]
    其余分量均为零。

    无迹约化四极矩的定义为：
    \[
        J_{ij} = q_{ij} - \frac{1}{3} \delta_{ij} \delta^{kl} q_{kl} = q_{ij} - \frac{1}{3} \delta_{ij} \Tr(q)
    \]
    因为$q_{xx} = q_{yy} = q_{zz}$，所以：
    \[
        J_{xx} = q_{xx} - \frac{1}{3} (q_{xx} + q_{yy} + q_{zz}) = q_{xx} - q_{xx} = 0
    \]
    同理，$J_{yy} = 0$，$J_{zz} = 0$，而其余非对角线分量也为零。因此对于一个球对称系统：
    \[
        J_{ij} = 0
    \]

    \item[8.(b)]
    对于一个具有均匀密度$\rho$的椭球体，椭球体内的$T^{00} = \rho$，则有：
    \[
        q_{xx} = \iiint_V x^2 \rho \dd^3{r} = \int_{-a}^a x^2\rho \dd{x} \iint_{D_x} \dd{y} \dd{z}
    \]
    其中$D_x$为平面$x = \mathrm{const.}$上椭圆的区域：
    \[
        \frac{y^2}{b^2} + \frac{z^2}{c^2} = 1 - \frac{x^2}{a^2}
    \]
    $D_x$是一个椭圆，其面积可以直接计算：
    \[
        S_x = \iint_{D_x} \dd{y} \dd{z} = \pi b c \qty(1 - \frac{x^2}{a^2})
    \]
    因此：
    \[
        q_{xx} = \int_{-a}^a x^2 \rho \pi b c \qty(1 - \frac{x^2}{a^2}) \dd{x} = \pi b c \rho \int_{-a}^a \qty(x^2 - \frac{x^4}{a^2}) \dd{x} = \frac{4\pi abc}{15} \rho a^2
    \]
    同理可得：
    \[
        q_{yy} = \frac{4\pi abc}{15} \rho b^2,\quad q_{zz} = \frac{4\pi abc}{15} \rho c^2
    \]
    其余分量均为零（因为被积函数是奇函数）。

    对应的无迹约化四极矩为：
    \[
        J_{xx} = q_{xx} - \frac{1}{3} (q_{xx} + q_{yy} + q_{zz}) = \frac{4\pi abc}{45} \rho \qty(2a^2 - b^2 - c^2)
    \]
    \[
        J_{yy} = q_{yy} - \frac{1}{3} (q_{xx} + q_{yy} + q_{zz}) = \frac{4\pi abc}{45} \rho \qty(2b^2 - a^2 - c^2)
    \]
    \[
        J_{zz} = q_{zz} - \frac{1}{3} (q_{xx} + q_{yy} + q_{zz}) = \frac{4\pi abc}{45} \rho \qty(2c^2 - a^2 - b^2)
    \]
    其余分量均为零。

    \item[8.(c)]
    对于点质量，可以用求和代替积分，设每个质量块质量为$m$，则有：
    \[
        T^{00} = \sum_{a=1}^{n} m \delta(\mathbf{r} - \mathbf{r}_a)
    \]
    \[
        q_{ij} = \sum_{a=1}^{n} m r_i^a r_j^a
    \]
    \[
        J_{ij} = \sum_{a=1}^n m \qty[r_i^a r_j^a - \frac{1}{3} \delta_{ij} (x_a^2 + y_a^2 + z_a^2)]
    \]
    在本题中，四个质量块的位置分别为：
    \[
        \mathbf{r}_1 = (a, 0, 0),\quad \mathbf{r}_2 = (-a, 0, 0),\quad \mathbf{r}_3 = (0, a, 0),\quad \mathbf{r}_4 = (0, -a, 0)
    \]
    则有：
    \[
        q_{xx} = \sum_{a=1}^{4} m x_a^2 = ma^2 + ma^2 + 0 + 0 = 2ma^2
    \]
    \[
        q_{yy} = \sum_{a=1}^{4} m y_a^2 = 0 + 0 + ma^2 + ma^2 = 2ma^2
    \]
    \[
        q_{zz} = \sum_{a=1}^{4} m z_a^2 = 0 + 0 + 0 + 0 = 0
    \]
    其余分量均为零。

    则四极矩的迹为：
    \[
        \Tr(q) = q_{xx} + q_{yy} + q_{zz} = 4ma^2
    \]
    无迹约化四极矩为：
    \[
        J_{xx} = q_{xx} - \frac{1}{3} \Tr(q) = 2ma^2 - \frac{4ma^2}{3} = \frac{2ma^2}{3}
    \]
    \[
        J_{yy} = q_{yy} - \frac{1}{3} \Tr(q) = 2ma^2 - \frac{4ma^2}{3} = \frac{2ma^2}{3}
    \]
    \[
        J_{zz} = q_{zz} - \frac{1}{3} \Tr(q) = 0 - \frac{4ma^2}{3} = -\frac{4ma^2}{3}
    \]
    其余分量均为零。

    \item[8.(d)]
    由题意，四个质点的坐标为：
    \[
        \mathbf{r}_1 = (a\cos{\omega t}, a\sin{\omega t}, 0),\quad \mathbf{r}_2 = (-a\sin{\omega t}, a\cos{\omega t}, 0)
    \]
    \[
        \mathbf{r}_3 = (-a\cos{\omega t}, -a\sin{\omega t}, 0),\quad \mathbf{r}_4 = (a\sin{\omega t}, -a\cos{\omega t}, 0)
    \]
    则质量四极矩为：
    \[
        q_{xx} = \sum_{a=1}^{4} m x_a^2 = m(2a^2 \cos^2{\omega t} + 2a^2 \sin^2{\omega t}) = 2ma^2
    \]
    \[
        q_{yy} = \sum_{a=1}^{4} m y_a^2 = m(2a^2 \sin^2{\omega t} + 2a^2 \cos^2{\omega t}) = 2ma^2
    \]
    \[
        q_{zz} = \sum_{a=1}^{4} m z_a^2 = 0
    \]
    其余分量均为零。

    则四极矩的迹为：
    \[
        \Tr(q) = q_{xx} + q_{yy} + q_{zz} = 4ma^2
    \]
    无迹约化四极矩为：
    \[
        J_{xx} = q_{xx} - \frac{1}{3} \Tr(q) = 2ma^2 - \frac{4ma^2}{3} = \frac{2ma^2}{3}
    \]
    \[
        J_{yy} = q_{yy} - \frac{1}{3} \Tr(q) = 2ma^2 - \frac{4ma^2}{3} = \frac{2ma^2}{3}
    \]
    \[
        J_{zz} = q_{zz} - \frac{1}{3} \Tr(q) = 0 - \frac{4ma^2}{3} = -\frac{4ma^2}{3}
    \]
    其余分量均为零。

    \item[8.(e)]
    由题意，弹簧的圆频率为：
    \[
        \omega = \sqrt{\frac{k(m + M)}{mM}}
    \]
    则两个质点的坐标为：
    \[
        \mathbf{r}_1 = \qty(b_0 \frac{M}{m + M} + \frac{2AM}{m + M}\cos{\omega t}, 0, 0),\quad \mathbf{r}_2 = \qty(-b_0 \frac{m}{m + M} - \frac{2Am}{m + M}\cos{\omega t}, 0, 0)
    \]
    则质量四极矩为：
    \begin{align*}
        q_{xx}
        &= m x_m^2 + M x_M^2 \\
        &= m\qty(b_0 \frac{M}{m + M} + \frac{2AM}{m + M}\cos{\omega t})^2 + M\qty(-b_0 \frac{m}{m + M} - \frac{2Am}{m + M}\cos{\omega t})^2 \\
        &= m \frac{M^2}{(m + M)^2} (b_0 + 2A\cos{\omega t})^2 + M \frac{m^2}{(m + M)^2} (-b_0 - 2A\cos{\omega t})^2 \\
        &= \frac{mM}{m + M} (b_0 + 2A\cos{\omega t})^2
    \end{align*}
    其余分量均为零。

    则四极矩的迹为：
    \[
        \Tr(q) = q_{xx} = \frac{mM}{m + M} (b_0 + 2A\cos{\omega t})^2
    \]
    无迹约化四极矩为：
    \[
        J_{xx} = q_{xx} - \frac{1}{3} \Tr(q) = \frac{2mM}{3(m + M)} (b_0 + 2A\cos{\omega t})^2
    \]
    \[
        J_{yy} = q_{yy} - \frac{1}{3} \Tr(q) = -\frac{mM}{3(m + M)} (b_0 + 2A\cos{\omega t})^2
    \]
    \[
        J_{zz} = q_{zz} - \frac{1}{3} \Tr(q) = -\frac{mM}{3(m + M)} (b_0 + 2A\cos{\omega t})^2
    \]
    其余分量均为零。

    \item[9.]
    设两颗星间距为$r$，并记$M = m_1 + m_2$，则轨道圆频率为：
    \[
        \omega = \sqrt{\frac{G(m_1 + m_2)}{r^3}} = \sqrt{\frac{GM}{r^3}}
    \]
    对于恒星1，其坐标为：
    \[
        \mathbf{r}_1 = \qty(\frac{m_2}{M} r \cos{\omega t}, \frac{m_2}{M} r \sin{\omega t}, 0)
    \]
    对于恒星2，其坐标为：
    \[
        \mathbf{r}_2 = \qty(-\frac{m_1}{M} r \cos{\omega t}, -\frac{m_1}{M} r \sin{\omega t}, 0)
    \]
    则质量四极矩为：
    \[
        q_{xx} = \sum_{i=1}^2 m_i x_i^2 = m_1 \qty(\frac{m_2}{M} r \cos{\omega t})^2 + m_2 \qty(-\frac{m_1}{M} r \cos{\omega t})^2 = \frac{m_1 m_2}{M} r^2 \cos^2{\omega t}
    \]
    \[
        q_{yy} = \sum_{i=1}^2 m_i y_i^2 = m_1 \qty(\frac{m_2}{M} r \sin{\omega t})^2 + m_2 \qty(-\frac{m_1}{M} r \sin{\omega t})^2 = \frac{m_1 m_2}{M} r^2 \sin^2{\omega t}
    \]
    此时有不为零且随时间变化的非对角分量：
    \begin{align*}
        q_{xy} = q_{yx}
        &= \sum_{i=1}^2 m_i x_i y_i \\
        &= m_1 \qty(\frac{m_2}{M} r \cos{\omega t}) \qty(\frac{m_2}{M} r \sin{\omega t}) + m_2 \qty(-\frac{m_1}{M} r \cos{\omega t}) \qty(-\frac{m_1}{M} r \sin{\omega t}) \\
        &= \frac{m_1 m_2^2 + m_2 m_1^2}{M^2} r^2 \cos{\omega t} \sin{\omega t} \\
        &= \frac{m_1 m_2}{M} r^2 \cos{\omega t} \sin{\omega t} \\
        &= \frac{m_1 m_2}{2M} r^2 \sin{2\omega t}
    \end{align*}
    其余分量均为零。

    则四极矩的迹为：
    \[
        \Tr(q) = q_{xx} + q_{yy} = \frac{m_1 m_2}{M} r^2 (\cos^2{\omega t} + \sin^2{\omega t}) = \frac{m_1 m_2}{M} r^2
    \]
    无迹约化四极矩为：
    \[
        J_{xx} = q_{xx} - \frac{1}{3} \Tr(q) = \frac{m_1 m_2}{M} r^2 \qty(\cos^2{\omega t} - \frac{1}{3})
    \]
    \[
        J_{yy} = q_{yy} - \frac{1}{3} \Tr(q) = \frac{m_1 m_2}{M} r^2 \qty(\sin^2{\omega t} - \frac{1}{3})
    \]
    \[
        J_{zz} = q_{zz} - \frac{1}{3} \Tr(q) = \frac{2m_1 m_2}{3M} r^2
    \]
    \[
        J_{xy} = J_{yx} = q_{xy} = \frac{m_1 m_2}{2M} r^2 \sin{2\omega t}
    \]
    其余分量均为零。

    之后计算$J_{ij}$对时间的三阶导数：
    \[
        \dddot{J}_{xx} = \frac{m_1 m_2}{M} r^2 \cdot 4\omega^3 \sin{2\omega t}
    \]
    \[
        \dddot{J}_{yy} = -\frac{m_1 m_2}{M} r^2 \cdot 4\omega^3 \sin{2\omega t}
    \]
    \[
        \dddot{J}_{xy} = \dddot{J}_{yx} = \frac{m_1 m_2}{M} r^2 \cdot 4\omega^3 \cos{2\omega t}
    \]
    其余分量均为零。

    根据引力辐射的辐射功率公式：
    \[
        P = \frac{G}{5} \left\langle\dddot{J}_{ij} \dddot{J}^{\,ij}\right\rangle
    \]
    则有：
    \begin{align*}
        P
        &= \frac{G}{5} \left\langle\dddot{J}_{ij} \dddot{J}^{\,ij}\right\rangle \\
        &= \frac{G}{5} \qty(\dddot{J}_{xx}^{\,2} + \dddot{J}_{yy}^{\,2} + \dddot{J}_{xy}^{\,2} + \dddot{J}_{yx}^{\,2}) \\
        &= \frac{G}{5} \qty(\frac{m_1 m_2}{M})^2 r^4 \cdot 16\omega^6 \qty(\sin^2{2\omega t} + \sin^2{2\omega t} + \cos^2{2\omega t} + \cos^2{2\omega t}) \\
        &= \frac{32G}{5} \qty(\frac{m_1 m_2}{M})^2 r^4 \omega^6
    \end{align*}
    再代入$M$和$\omega$的表达式，辐射功率为：
    \[
        P = \frac{32G^4}{5} \frac{(m_1 m_2)^2 (m_1 + m_2)}{r^5}
    \]
    为了求出两颗恒星的轨道运动圆频率$\omega$随时间的变化率$\dot{\omega}$，利用轨道总能量：
    \[
        E = -\frac{Gm_1 m_2}{2r} \implies \dot{E} = -P = \frac{G m_1 m_2}{2r^2} \dot{r}
    \]
    则有：
    \[
        \dot{r} = -\frac{2r^2}{G m_1 m_2} P = -\frac{64G^3}{5} \frac{m_1 m_2 (m_1 + m_2)}{r^3}
    \]
    根据$\omega$与$r$的关系：
    \[
        \omega = \sqrt{\frac{G(m_1 + m_2)}{r^3}}
    \]
    则有：
    \[
        \dot{\omega} = -\frac{3}{2} \sqrt{\frac{G(m_1 + m_2)}{r^5}} \dot{r}
    \]
    代入$\dot{r}$的表达式，得到：
    \[
        \dot{\omega} = \frac{96G^{5/2}}{5} \frac{(m_1 m_2) \sqrt{m_1 + m_2}}{r^{11/2}}
    \]

\end{enumerate}

\end{document}